\section{Background}
\label{sec:background}

% (~8-10 pages)

\subsection{The Human Gut Microbiome}
% Microbial ecology of the gut
% Key phyla: Firmicutes, Bacteroidetes, Proteobacteria, Actinobacteria
% Functional roles: nutrient metabolism, immune training, pathogen resistance
% The Firmicutes/Bacteroidetes ratio as a crude health indicator
% Why composition alone is insufficient — need structural understanding

\subsection{Microbiome Dysbiosis and the Permeability Cascade}
\label{subsec:permeability_cascade}
% Definition: loss of beneficial organisms, expansion of pathobionts,
%   reduction in overall diversity
% The cascade mechanism:
%   1. Dysbiosis → reduced butyrate-producing bacteria (Faecalibacterium,
%      Roseburia, Eubacterium)
%   2. Butyrate depletion → reduced tight junction protein expression
%      (claudins, occludin, ZO-1)
%   3. Barrier degradation → increased intestinal permeability ("leaky gut")
%   4. LPS and other endotoxins translocate to systemic circulation
%   5. TLR4 activation → NF-κB → pro-inflammatory cytokines
%      (TNF-α, IL-1β, IL-6)
%   6. Systemic inflammation → neuroinflammation via blood-brain barrier
%      and vagal afferents
% Permeability biomarkers: zonulin, LPS-binding protein (LBP),
%   intestinal fatty acid-binding protein (I-FABP),
%   lactulose/mannitol ratio
% Associated diseases: IBD, IBS, obesity, T2D, depression, anxiety,
%   Parkinson's, Alzheimer's

\subsection{The Gut-Brain Axis}
\label{subsec:gut_brain}
% Bidirectional communication via:
%   - Vagus nerve (primary conduit)
%   - Enteric nervous system (ENS, "second brain")
%   - Immune signaling (cytokines crossing BBB)
%   - Microbial metabolites (SCFAs, tryptophan metabolites)
% Neuropeptides modulated by gut microbiota:
%   - Serotonin (5-HT): 95\% produced by enterochromaffin cells in gut,
%     microbiota regulate tryptophan hydroxylase expression
%   - Ghrelin: appetite regulation, modulated by gut composition
%   - Cholecystokinin (CCK): satiety, gut motility
%   - Vasoactive intestinal peptide (VIP): anti-inflammatory,
%     neuroprotective, barrier maintenance
%   - GABA: produced by Lactobacillus and Bifidobacterium spp.
% Disrupted signaling in dysbiosis → cognitive impairment,
%   emotional dysregulation, altered decision-making

\subsection{Mycotoxins and Environmental Triggers of Dysbiosis}
\label{subsec:mycotoxins}
% Mycotoxins as a concrete environmental trigger
% Trichothecenes (T-2 toxin, deoxynivalenol/DON):
%   - Directly reduce tight junction protein expression
%   - Decrease mucin-secreting goblet cells
%   - Alter microbiome composition (reduce Lactobacillus, increase
%     Proteobacteria)
% Aflatoxins, ochratoxin A: additional gut barrier effects
% Chronic Inflammatory Response Syndrome (CIRS):
%   - Biotoxin illness from water-damaged buildings
%   - Involves gut barrier dysfunction, immune dysregulation,
%     neuroinflammation
%   - HLA-DR genetic susceptibility (~24\% of population)
%   - Shoemaker protocol: cholestyramine, VIP nasal spray
% Mycotoxin-induced dysbiosis → bacterial translocation →
%   chronic inflammation → neurological effects
% This provides a concrete use case for topological regime detection

\subsection{Co-occurrence Networks in Microbial Ecology}
\label{subsec:cooccurrence}
% Why networks: taxa don't exist in isolation
% Network construction from abundance data:
%   - Pearson/Spearman correlation (problematic for compositional data)
%   - SparCC: handles compositionality via log-ratio variance
%   - SPIEC-EASI: sparse inverse covariance estimation
% Network metrics: modularity, clustering coefficient, hub scores,
%   betweenness centrality
% Ecological interpretation: modules = functional guilds,
%   hubs = keystone species
% Limitations: pairwise relationships miss higher-order structure

\subsection{Topological Data Analysis}
\label{subsec:tda}
% Simplicial complexes: generalizing graphs to higher dimensions
%   - 0-simplices (vertices), 1-simplices (edges),
%     2-simplices (triangles), etc.
% Filtrations: growing simplicial complexes by a parameter
%   - Vietoris-Rips: add simplex when all pairwise distances < ε
%   - \v{C}ech: add simplex when balls of radius ε have common intersection
% Persistent homology: tracking topological features across filtration
%   - Birth and death of features
%   - Persistence = death - birth (long-lived features are "real")
% Betti numbers and their ecological interpretation:
%   - β₀: connected components (isolated community clusters)
%   - β₁: loops (ecological feedback cycles, mutualism rings)
%   - β₂: voids (higher-order community cavities)
% Persistence diagrams, barcodes, landscapes
% Prior applications:
%   - Protein structure and function \citep{camara2017topological}
%   - Neural spike train analysis
%   - Financial market regime detection (the parallel to this work)
%   - Biological aggregation models \citep{topaz2015topological}
%   - Zebrafish pattern formation \citep{mcguirl2020topological}
